\section{课后测验题目题解}

\subsection{第一章测试题解}

\subsubsection{单选题}

\begin{enumerate}
	\item \textbf{解：}
	
	对于集合 $P$：\\
	由 $(x+1)^2 \le 4$，两边开平方得：
	\[ -2 \le x+1 \le 2 \]
	解得：
	\[ -3 \le x \le 1 \]
	所以，集合 $P = [-3, 1]$。
	
	对于集合 $Q$：\\
	由 $5 \le x^2 + 16$，移项化简得：
	\[ x^2 + 11 \ge 0 \]
	由于 $x \in \mathbb{R}$，实数平方非负（$x^2 \ge 0$），故 $x^2 + 11 \ge 11 > 0$ 恒成立。\\
	所以，集合 $Q = \mathbb{R}$。
	
	对比两集合的范围：\\
	因为 $P = [-3, 1]$ 且 $Q = \mathbb{R}$，显然集合 $P$ 中的所有元素均属于集合 $Q$，且 $Q$ 中存在不属于 $P$ 的元素。即 $P$ 是 $Q$ 的真子集：
	\[ P \subset Q \]
	
	故选 D。
	
	\item \textbf{解：}
	
	已知集合 $A = \{a, \{a\}\}$，该集合包含两个元素：$a$ 和 $\{a\}$。\\
	幂集 $P(A)$ 是由集合 $A$ 的所有子集构成的集合。因为 $A$ 含有 2 个元素，所以 $P(A)$ 共有 $2^2 = 4$ 个元素，具体枚举如下：
	\[ P(A) = \{\emptyset, \{a\}, \{\{a\}\}, \{a, \{a\}\}\} \]
	
	逐项分析：\\
	\textbf{A 选项}：$\{a\}$ 是 $P(A)$ 中的一个元素，故 $\{a\} \in P(A)$ 命题正确。\\
	\textbf{B 选项}：$\{\{a\}\}$ 是 $P(A)$ 中的一个元素，故 $\{\{a\}\} \in P(A)$ 命题正确。\\
	\textbf{C 选项}：若 $\{a\} \subseteq P(A)$ 成立，根据子集的定义，必有 $a \in P(A)$。但观察幂集 $P(A)$ 的元素可知，$a$ 并非 $P(A)$ 的元素（即 $a \notin P(A)$），故该命题错误。\\
	\textbf{D 选项}：若 $\{\{a\}\} \subseteq P(A)$ 成立，则必有 $\{a\} \in P(A)$。根据 A 选项的结论，这显然成立，故该命题正确。
	
	题目要求选出错误的命题，因此正确答案是 \textbf{C}。
	
	\item \textbf{解：}
	
	已知 $0$ 是一个元素，而 $\emptyset$ 表示空集。\\
	描述元素与集合的关系用属于（$\in$）或不属于（$\notin$）。由于空集中不包含任何元素，故 $0$ 不属于空集，即：
	\[ 0 \notin \emptyset \]正确符号为 $\notin$\\
	故选 \textbf{D}。
	
	\item \textbf{解:}\\
	已知 $A - B = \emptyset$（原题中 $\Phi$ 视为空集）。\\
	差集 $A - B$ 表示属于 $A$ 且不属于 $B$ 的元素集合。$A - B = \emptyset$ 意味着 $A$ 中不存在任何不属于 $B$ 的元素，即 $A$ 的所有元素均属于 $B$，推导出 $A \subseteq B$。
	逐项分析：\\
		\textbf{A 选项}：$A \subset B$ 是可能的（如 $A=\{1\}, B=\{1, 2\}$）。\\
		\textbf{B 选项}：$A = \emptyset$ 是可能的（此时必定有 $A-B=\emptyset$）。\\
		\textbf{C 选项}：$B \subset A$ 表示 $B$ 是 $A$ 的真子集，这意味着 $A$ 中存在不属于 $B$ 的元素，此时必然有 $A - B \neq \emptyset$，与已知矛盾，故\textbf{不可能正确}。\\
		\textbf{D 选项}：$B = \emptyset$ 是可能的（当 $A=\emptyset$ 时成立）。\\

	故选 \textbf{C}。
	
	\item \textbf{解:}\\
	逐项分析命题的真假：\\
		\textbf{A 选项}：空集是任何集合的子集，包括空集自身（$\emptyset \subseteq \emptyset$），所以不仅仅是非空集合的子集。命题错误。\\
		\textbf{B 选项}：仅凭 $A$ 的一个元素属于 $B$ 无法推断两集合相等（如 $A=\{1, 2\}, B=\{1, 3\}$）。命题错误。\\
		\textbf{C 选项}：若 $A \neq B$，则必存在元素只属于其中一个集合。例如存在 $x \in A$ 且 $x \notin B$，则 $x \in A-B$ 且 $x \notin B-A$，导致 $A-B \neq B-A$。故若要 $A-B = B-A$ 成立，必须没有任何元素单独属于某一个集合，即 $A = B$。命题正确。\\
		\textbf{D选项}：空集是任何\textbf{非空}集合的真子集。空集不是自身的真子集（$\emptyset \not\subset \emptyset$）。命题错误。\\
	故选 \textbf{C}。
	
	\item \textbf{解：}\\
	题目要求寻找代数系统 $\langle \{a, b, c\}, * \rangle$ 的零元。\\
	\textbf{零元的定义是：}对于集合中的任意元素 $x$，均满足 $x * z = z * x = z$，则称 $z$ 为零元。\\
	观察运算表：\\
	对于元素 $c$，有 $a * c = c$，$b * c = c$，$c * c = c$（即第 $c$ 列全为 $c$）；\\
	且 $c * a = c$，$c * b = c$，$c * c = c$（即第 $c$ 行全为 $c$）。\\
	这说明对于任意 $x \in \{a, b, c\}$，都有 $x * c = c * x = c$ 成立。\\
	因此，$c$ 是该代数系统的零元。\\
	故选 \textbf{D}。
	
	\item \textbf{第 7 题：}\\
	已知函数 $f: \mathbb{N} \to \mathbb{N}$，定义为 $f(x) = x \bmod 3$。\\
	判断单射性：\\
	取 $x_1 = 0, x_2 = 3$（$0, 3 \in \mathbb{N}$），计算得 $f(0) = 0 \bmod 3 = 0$，$f(3) = 3 \bmod 3 = 0$。\\
	因为 $f(0) = f(3)$ 但 $0 \neq 3$，所以 $f$ \textbf{不是单射}。\\
	判断满射性：\\
	任意自然数除以 $3$ 的余数只能是 $0, 1, 2$。因此函数 $f$ 的值域为 $\{0, 1, 2\}$。\\
	由于值域 $\{0, 1, 2\}$ 严格包含于陪域 $\mathbb{N}$（例如 $3 \in \mathbb{N}$，但不存在 $x \in \mathbb{N}$ 使 $x \bmod 3 = 3$），所以 $f$ \textbf{不是满射}。\\
	综上，$f$ 既不是单射也不是满射。\\
	故选 \textbf{D}。
	
	\item \textbf{解:}\\
	题目要求选出是单射而非满射的函数。逐项分析如下：\\
	\textbf{A 选项}：$f: \mathbb{R} \to \mathbb{R}, f(x) = 2x + 1$。由 $2x_1+1 = 2x_2+1 \Rightarrow x_1=x_2$ 知其为单射；对任意 $y \in \mathbb{R}$，存在 $x = (y-1)/2 \in \mathbb{R}$ 使 $f(x)=y$，知其为满射。该函数是双射，不符合题意。\\
	\textbf{B 选项}：$f: \mathbb{R} \to \mathbb{Z}, f(x) = \lfloor x \rfloor$。由于 $f(1.1) = 1$ 且 $f(1.2) = 1$，不同自变量对应相同函数值，故不是单射。不符合题意。\\
	\textbf{C 选项}：$f: \mathbb{R} \to \mathbb{R}, f(x) = -x^2 + 2x - 1 = -(x-1)^2$。由于 $f(0) = -1$ 且 $f(2) = -1$，故不是单射。不符合题意。\\
	\textbf{D 选项}：$f: \mathbb{Z}^+ \to \mathbb{R}, f(x) = \ln x$。因为对数函数在定义域上严格单调递增，故 $x_1 \neq x_2$ 时必有 $\ln x_1 \neq \ln x_2$，它是单射；而函数的值域为 $\{0, \ln 2, \ln 3, \dots\}$，这显然只是实数集 $\mathbb{R}$ 的一个真子集（例如实数 $0.5$ 就不在值域内，因为解得 $x = e^{0.5} \notin \mathbb{Z}^+$），故不是满射。符合题意。\\
	故选 \textbf{D}。
\end{enumerate}

\subsubsection{填空题}

\begin{enumerate}
	\item \textbf{解：}
	
	若 $S = \{S_1, S_2, \cdots, S_m\}$ 是集合 $A$ 的一个分划（划分），则它必须同时满足以下三个条件：\\
	(1) 划分中的所有子集非空：$S_i \neq \emptyset$ （对于 $1 \le i \le m$）；\\
	(2) 划分中的任意两个不同子集互不相交：当 $i \neq j$ 时，$S_i \cap S_j = \emptyset$；\\
	(3) 所有子集的并集等于原集合：$\bigcup_{i=1}^m S_i = A$。
	
	\vspace{1em}
	
	\item \textbf{解:}\\
	求并集 $A \cup B$：将两集合中所有的有序对合并，并去除重复项，得到：\\
	$A \cup B = \{\langle 1,2 \rangle, \langle 2,4 \rangle, \langle 3,3 \rangle, \langle 1,3 \rangle, \langle 4,2 \rangle\}$。\\
	求复合关系 $A \circ B$：根据关系复合的左复合规则，若存在 $y$ 使得 $\langle x,y \rangle \in A$ 且 $\langle y,z \rangle \in B$，则 $\langle x,z \rangle \in A \circ B$。\\
	因为 $\langle 1,2 \rangle \in A$ 且 $\langle 2,4 \rangle \in B$，得 $\langle 1,4 \rangle \in A \circ B$；\\
	因为 $\langle 2,4 \rangle \in A$ 且 $\langle 4,2 \rangle \in B$，得 $\langle 2,2 \rangle \in A \circ B$；\\
	有序对 $\langle 3,3 \rangle \in A$ 的第二元为 3，而集合 $B$ 中没有以 3 为第一元的有序对，故无对应。\\
	所以 $A \circ B = \{ \langle 1,4 \rangle, \langle 2,2 \rangle \}$。
	
	\vspace{1em}
	
	\item \textbf{解：}\\
	对于任意有限集合 $S$，其幂集 $P(S)$ 的基数公式为 $|P(S)| = 2^{|S|}$。\\
	已知 $|S| = 5$，则 $|P(S)| = 2^5 = 32$。
	
	\vspace{1em}
	
	\item \textbf{解:}\\
	观察文氏图，阴影部分严格位于集合 $C$ 的内部。同时，它完全排除了与 $A$ 产生交集的区域，也完全排除了与 $B$ 产生交集的区域。这意味着该区域包含了所有属于 $C$ 但既不属于 $A$ 也不属于 $B$ 的元素。\\
	其集合表达式可写作：$(B \oplus C - A)$。
	
	\vspace{1em}
	
	\item \textbf{解:}\\
	题目中以 $\Phi$ 表示空集 $\emptyset$。逐一判断命题真假：\\
	\textbf{(1)}：集合 $\{\Phi, \{\{\Phi\}\}\}$ 仅包含两个元素：$\Phi$ 和 $\{\{\Phi\}\}$。元素 $\{\Phi\}$ 并不在其中，故 \textbf{(1) 错误}。\\
	\textbf{(2)}：判断 $\{\Phi\} \subset \{\Phi, \{\{\Phi\}\}\}$。左边集合的唯一元素是 $\Phi$，而 $\Phi$ 确实是右侧集合的元素，且右侧集合还有其他元素。满足真子集定义，故 \textbf{(2) 正确}。\\
	\textbf{(3)}：集合 $\{\{\Phi\}\}$ 的唯一元素是 $\{\Phi\}$。因为 $\Phi \neq \{\Phi\}$，所以 $\Phi$ 不属于该集合，故 \textbf{(3) 错误}。\\
	\textbf{(4)}：空集是任何集合的子集，故 $\Phi \subseteq \{\Phi\}$ 恒成立，故 \textbf{(4) 正确}。\\
	\textbf{(5)}：右侧集合 $\{a,b, \{b\}, \{a\}\}$ 的元素依次是 $a$, $b$, $\{b\}$, $\{a\}$。并没有把 $\{a,b\}$ 作为一个整体集合元素列出，故 \textbf{(5) 错误}。\\
	综上所述，正确的命题是 \textbf{(2)} 和 \textbf{(4)}。
	
	\vspace{1em}
	
	\item \textbf{解：}\\
	逐个分析推理过程：\\
	\textbf{(1)}：$A \subseteq B, B \subseteq C \Rightarrow A \subseteq C$。这是集合包含关系的传递性，显然成立，推理\textbf{正确}。\\
	\textbf{(2)}：$A \subseteq B, B \subseteq C \Rightarrow A \in B$。集合的包含关系（$\subseteq$）不能推导出元素的从属关系（$\in$）。例如 $A=\{1\}, B=\{1,2\}, C=\{1,2,3\}$，满足前置条件，但 $\{1\} \notin B$，推理\textbf{错误}。\\
	\textbf{(3)}：$A \in B, B \in C \Rightarrow A \in C$。元素的从属关系不具备传递性。例如 $A=1, B=\{1\}, C=\{\{1\}\}$，满足 $1 \in B$ 且 $B \in C$，但 $1 \notin C$，推理\textbf{错误}。\\
	综上，正确的推理只有 \textbf{(1)}。
	
	\vspace{1em}
	
	\item \textbf{解:}\\
	\textbf{寻找幺元（单位元）：}\\
	幺元 $e$ 需要满足对任意元素 $x$，都有 $x * e = e * x = x$。\\
	观察运算表，在第 $a$ 行中，$a * a = a, a * b = b, a * c = c, a * d = d$；在第 $a$ 列中，$a * a = a, b * a = b, c * a = c, d * a = d$。这说明元素 $a$ 与任何元素运算都等于该元素本身。\\
	所以，该代数系统的\textbf{幺元是 $a$}。
	
	\textbf{寻找逆元：}\\
	若元素 $x$ 有逆元 $y$，则需满足 $x * y = y * x = a$（因为幺元是 $a$）。\\
	对于 $a$：因为 $a * a = a$，所以 $a$ 的逆元是 $a$。\\
	对于 $b$：因为 $b * d = a$ 且 $d * b = a$，所以 $b$ 的逆元是 $d$。\\
	对于 $c$：因为 $c * c = a$，所以 $c$ 的逆元是 $c$。\\
	对于 $d$：因为 $d * b = a$ 且 $b * d = a$，所以 $d$ 的逆元是 $b$。\\
	因此，\textbf{有逆元的元素为 $a, b, c, d$，它们的逆元分别为 $a, d, c, b$}。
	
	\vspace{1em}
	
	\item \textbf{解:}\\
	分别化简各集合中的元素：\\
	\textbf{(1)} $A_1 = \{a, b\}$\\
	\textbf{(2)} $A_2 = \{b, a\}$。集合元素无序，故 $A_2 = \{a, b\}$\\
	\textbf{(3)} $A_3 = \{a, b, a\}$。集合元素具有互异性，去除重复元素后 $A_3 = \{a, b\}$\\
	\textbf{(4)} $A_4 = \{a, b, c\}$\\
	\textbf{(5)} $A_5 = \{x \mid (x-a)(x-b)(x-c) = 0\}$。该方程的解为 $x=a, x=b, x=c$，故该集合的元素为 $a, b, c$，即 $A_5 = \{a, b, c\}$\\
	\textbf{(6)} $A_6 = \{x \mid x^2-(a+b)x+ab = 0\}$。该一元二次方程可因式分解为 $(x-a)(x-b) = 0$，解为 $x=a, x=b$，故该集合的元素为 $a, b$，即 $A_6 = \{a, b\}$
	
	对比化简结果，相等的集合分为两组：\\
	\textbf{第一组：} $A_1 = A_2 = A_3 = A_6$\\
	\textbf{第二组：} $A_4 = A_5$
	
	\vspace{1em}
	
	\item \textbf{解：}\\
	已知 $N$ 为自然数集（在离散数学中通常包含 $0$），则由 $x \in N$ 且 $x < 5$ 可得集合 $A = \{0, 1, 2, 3, 4\}$。\\
	已知 $E^+$ 为正偶数集，由 $x \in E^+$ 且 $x < 7$ 可得集合 $B = \{2, 4, 6\}$。\\
	求并集 $A \cup B$，即将两集合元素合并并去除重复项，得到：
	\[ A \cup B = \{0, 1, 2, 3, 4, 6\} \]
	
	\vspace{1em}
	
	\item \textbf{解:}\\
	逐一判断命题真假：\\
	\textbf{(1)}：若 $A \cup B = A \cup C$，则 $B = C$。此命题\textbf{错误}。举个反例：设 $A=\{1\}, B=\{1\}, C=\emptyset$。此时 $A \cup B = \{1\}$ 且 $A \cup C = \{1\}$，满足前提条件，但显然 $B \neq C$。\\
	\textbf{(2)}：$\{a, b\} = \{b, a\}$。集合中的元素具有无序性，无论元素列出的先后顺序如何，集合都相等。此命题\textbf{正确}。\\
	\textbf{(3)}：$P(A \cap B) \neq P(A) \cap P(B)$。此命题\textbf{错误}。在集合论中，交集的幂集等于幂集的交集，即 $P(A \cap B) = P(A) \cap P(B)$ 是一个恒成立的定理。\\
	\textbf{(4)}：若 $A$ 为非空集，则 $A \neq A \cup A$ 成立。此命题\textbf{错误}。根据集合运算的等幂律，对于任意集合 $A$（无论是否为空），都有 $A \cup A = A$ 恒成立。\\
	综上所述，正确的命题只有 \textbf{(2)}。
	
	\vspace{1em}
	
	\item \textbf{解:}\\
	\textbf{寻找幺元：}\\
	幺元 $e$ 需要满足：对于集合 $S$ 中的任意元素 $x$，都有 $e * x = x * e = x$ 成立。\\
	观察给定的运算表，在第 $\beta$ 行中，$\beta * \alpha = \alpha, \beta * \beta = \beta, \beta * \gamma = \gamma, \beta * \delta = \delta$（运算结果与列首元素完全一致）；在第 $\beta$ 列中，$\alpha * \beta = \alpha, \beta * \beta = \beta, \gamma * \beta = \gamma, \delta * \beta = \delta$（运算结果与行首元素完全一致）。\\
	这说明元素 $\beta$ 与任何元素进行 $*$ 运算，结果都等于该元素本身。因此，该代数系统的\textbf{幺元是 $\beta$}。
	
	\textbf{寻找 $\alpha$ 的逆元：}\\
	已知幺元为 $\beta$，若 $\alpha$ 的逆元为 $x$，则需要满足 $\alpha * x = x * \alpha = \beta$。\\
	查表可知，在 $\alpha$ 所在的行中，$\alpha * \gamma = \beta$；在 $\alpha$ 所在的列中，$\gamma * \alpha = \beta$。\\
	满足条件的元素是 $\gamma$，因此，\textbf{$\alpha$ 的逆元是 $\gamma$}。
	
	\vspace{1em}
	
	\item \textbf{解:}\\
	题目要求计算复合函数 $f \circ g(x)$。\\
	根据复合函数的定义，$f \circ g(x)$ 等价于将 $g(x)$ 的值作为函数 $f$ 的自变量，即 $f(g(x))$。\\
	已知 $g(x) = 2x$，将其代入函数 $f$ 中：
	\[ f(g(x)) = f(2x) \]
	又因为对于任意 $x \in N$，都有 $f(x) = x + 1$，我们将自变量替换为 $2x$，得到：
	\[ f(2x) = 2x + 1 \]
	因此，$f \circ g(x) = 2x + 1$。
\end{enumerate}
\newpage

\subsection{第二章测试题解}

\subsubsection{选择题}

\begin{enumerate}
	\item \textbf{解:}
	
	观察给定的关系图，集合 $S=\{1, 2, 3\}$。图中的有向边代表关系 $R$ 中的元素。从图中可以看出，箭头构成了从 $1 \to 2$，$2 \to 3$，$3 \to 1$ 的有向环。因此，关系集合可以表示为：$R = \{(1,2), (2,3), (3,1)\}$。
	
	下面根据二元关系的性质对选项进行逐一分析：
	
	\textbf{A 选项}：自反性。自反性要求对于所有 $x \in S$，都有 $(x,x) \in R$。图中没有任何节点有指向自身的自环（即不存在 $(1,1), (2,2), (3,3)$），所以 $R$ 不具有自反性。此选项错误。
	
	\textbf{B 选项}：反自反性、反对称性、传递性。反自反性要求对于所有 $x \in S$，都有 $(x,x) \notin R$，由于图中无自环，所以 $R$ 具有反自反性。但 $R$ 不具有传递性，因为传递性要求若 $(x,y) \in R$ 且 $(y,z) \in R$，则必然 $(x,z) \in R$。在这里 $(1,2)\in R$ 且 $(2,3)\in R$，但是 $(1,3)\notin R$（图中只有 $3 \to 1$ 的箭头，无 $1 \to 3$ 的箭头）。此选项错误。
	
	\textbf{C 选项}：反自反性、反对称性。前面已证具有反自反性。反对称性要求如果 $(x,y) \in R$ 且 $(y,x) \in R$，则必然 $x=y$。也就是说，不同元素之间不能存在双向箭头。图中 $R$ 没有任何两个节点之间存在双向箭头，因此反对称性的条件前件始终为假，命题恒真，所以 $R$ 具有反对称性。此选项正确。
	
	\textbf{D 选项}：自反性、对称性、传递性。前面已证 $R$ 既不自反也不传递。同时 $R$ 也不对称（例如 $(1,2) \in R$ 但 $(2,1) \notin R$）。此选项错误。
	
	\textbf{正确答案：C}
	%\hrule
	\vspace{1em}
	
	\item \textbf{解:}
	
	题目问：设 $|A| = n$，则 $A$ 上有多少个二元关系。
	
	根据定义，集合 $A$ 上的二元关系 $R$ 是笛卡尔积 $A \times A$ 的任意一个子集。首先，计算笛卡尔积 $A \times A$ 中元素的总数：$|A \times A| = |A| \times |A| = n \times n = n^2$。其次，一个包含 $n^2$ 个元素的集合，其所有可能子集（即幂集）的个数为 $2^{n^2}$ 个。因此，集合 $A$ 上的二元关系总共有 $2^{n^2}$ 个。
	
	\textbf{A 选项}：$2^n$ 是集合 $A$ 的子集个数，也就是 $A$ 的一元关系的个数，而不是二元关系的个数。此选项错误。
	
	\textbf{B 选项}：$2^{n^2}$ 正好等于笛卡尔积 $A \times A$ 的子集总数，即二元关系的个数。此选项正确。
	
	\textbf{C 选项}：$n^2$ 是笛卡尔积 $A \times A$ 中的元素个数，代表的是全域关系中包含的有序对数量，而不是关系的种类数。此选项错误。
	
	\textbf{D 选项}：$n^n$ 是从集合 $A$ 映射到集合 $A$ 的所有可能函数的个数。由于每一个函数也是一种特殊的二元关系，所以 $n^n$ 只是关系的一部分，并不是所有二元关系的总数。此选项错误。
	
	\textbf{正确答案：B}

	\item \textbf{解：}\\
	已知集合 $A = \{\emptyset, \{1\}, \{1, 3\}, \{1, 2, 3\}\}$。\\
	分析集合 $A$ 中各个元素之间的包含关系（$\subseteq$）：\\
	显然有 $\emptyset \subseteq \{1\}$；\\
	$\{1\} \subseteq \{1, 3\}$；\\
	$\{1, 3\} \subseteq \{1, 2, 3\}$。\\
	由于集合中的包含关系具有传递性，这四个元素实际上构成了一个全序关系（即任意两个元素都是可比的，形成一条单向的包含链）。\\
	在绘制哈斯图（Hasse Diagram）时，包含关系表示为自下而上的连线，且省略自环和传递边。因此，最底层应当是最小元 $\emptyset$，往上依次用线段连接 $\{1\}$，然后是 $\{1, 3\}$，最顶层是最大元 $\{1, 2, 3\}$。整体图形必须是一条垂直的直线。\\
	观察四个选项，只有 A 选项的图形是一条垂直的链条。\\
	故选 \textbf{A}。
	
	\vspace{1em}
	
	\item \textbf{解:}\\
	题目要求判断哪个关系能构成函数。在集合论中，一个二元关系若要构成函数，必须满足“右唯一性”（单值性）：即对于定义域中的每一个输入值 $x_1$，有且仅有一个唯一的输出值 $x_2$ 与之对应。\\
	逐项进行分析：\\
	\textbf{A 选项}：$f = \{\langle x_1, x_2 \rangle \mid x_1, x_2 \in \mathbb{R}, x_1 = x_2^2\}$。假设输入 $x_1 = 4$，则满足 $4 = x_2^2$ 的 $x_2$ 有 $2$ 和 $-2$ 两个值。即 $\langle 4, 2 \rangle \in f$ 且 $\langle 4, -2 \rangle \in f$。一个输入对应了多个输出，不满足单值性，不能构成函数。\\
	\textbf{B 选项}：$f = \{\langle x_1, x_2 \rangle \mid (x_1, x_2 \in \mathbb{N}) \land (x_1 + x_2 = 10)\}$。题目通常默认函数是指在指定集合上的全函数。对于定义域 $\mathbb{N}$ 中的元素（例如取 $x_1 = 11$），解得 $x_2 = -1$，但 $-1 \notin \mathbb{N}$。这意味着定义域中存在部分元素找不到对应的映射值，因此它不是 $\mathbb{N}$ 上的全函数。\\
	\textbf{C 选项}：$f = \{\langle x, |x| \rangle \mid x \in \mathbb{R}\}$。在这个关系中，第一元是 $x$，第二元是 $|x|$。对于实数集 $\mathbb{R}$ 中的任意一个实数 $x$，它的绝对值 $|x|$ 都是唯一确定的实数。对于每一个输入，都有且仅有一个确定的输出，完全符合函数的定义（这就是标准的绝对值函数）。\\
	\textbf{D 选项}：$f = \{\langle x_1, x_2 \rangle \mid x_1, x_2 \in \mathbb{N}, x_2 \text{为小于} x_1 \text{的素数的个数}\}$。若自然数集 $\mathbb{N}$ 的定义不包含 $0$（部分国内旧教材标准 $\mathbb{N}=\{1,2,3...\}$），当 $x_1=1$ 时，小于 1 的素数个数为 0，而 $0 \notin \mathbb{N}$，此时找不到映射值。相较之下，C 选项是数学中最基础、无任何定义域歧义的函数形式。\\
	故选 \textbf{C}。
	
	\vspace{1em}
	
	\item \textbf{解：}\\
	已知集合 $S = \{1, 2, 3\}$。观察关系 $R$ 的关系图，图中的三个顶点 1, 2, 3 各有一个自环，且不存在任何其他的有向边（即不同顶点间没有边）。这说明关系 $R$ 的集合表示为：
	\[ R = \{\langle 1, 1 \rangle, \langle 2, 2 \rangle, \langle 3, 3 \rangle\} \]
	这就是集合 $S$ 上的恒等关系。我们逐一分析其性质：\\
	\textbf{1. 自反性：} 根据定义，对任意 $x \in S$，若都有 $\langle x, x \rangle \in R$，则关系自反。由于 $S$ 中的 1, 2, 3 均在 $R$ 中有对应的有序对（都有自环），故自反性成立。\\
	\textbf{2. 对称性：} 根据定义，若 $\langle x, y \rangle \in R$，则必有 $\langle y, x \rangle \in R$。在关系 $R$ 中，所有的元素均为 $\langle x, x \rangle$ 的形式，交换 $x$ 和 $y$ 的位置后仍为 $\langle x, x \rangle \in R$，故条件成立（满足“空真”），对称性成立。\\
	\textbf{3. 反对称性：} 根据定义，若 $\langle x, y \rangle \in R$ 且 $\langle y, x \rangle \in R$，则必有 $x = y$。因为 $R$ 中不存在 $x \neq y$ 的有序对，前提条件中的冲突不可能发生，故条件成立（满足“空真”），反对称性成立。\\
	\textbf{4. 传递性：} 根据定义，若 $\langle x, y \rangle \in R$ 且 $\langle y, z \rangle \in R$，则必有 $\langle x, z \rangle \in R$。由于 $R$ 中仅有恒等对，即 $x=y$ 且 $y=z$，必然推导出 $x=z$，即 $\langle x, z \rangle \in R$ 恒成立，故传递性成立。\\
	综上所述，$R$ 具有自反、对称、反对称、传递的性质。\\
	故选 \textbf{D}。
	
	\vspace{1em}
	
	\item \textbf{解:}\\
	已知集合 $A = \{1, 2, 3\}$，其包含的元素个数（即基数）为 $|A| = 3$。\\
	根据定义，集合 $A$ 上的二元关系是笛卡尔积 $A \times A$ 的子集。\\
	首先计算笛卡尔积 $A \times A$ 的元素个数：
	\[ |A \times A| = |A| \times |A| = 3 \times 3 = 3^2 \]
	所以 $A \times A$ 中共有 $3^2$ 个不同的有序对。\\
	再根据集合论原理，对于一个含有 $n$ 个元素的集合，其所有可能子集（即幂集）的数量为 $2^n$ 个。\\
	将 $n = 3^2$ 代入，可知 $A \times A$ 的子集总数为 $2^{3^2}$ 个，这就代表了集合 $A$ 上所有可能的二元关系的总数。\\
	故选 \textbf{A}。
	
	\vspace{1em}
	
	\item \textbf{解：}\\
	已知 $R, S$ 是集合 $A$ 上的关系。考察选项 D，若 $R, S$ 都是自反的，则对于任意 $x \in A$，均有 $(x, x) \in R$ 且 $(x, x) \in S$。根据复合关系 $R \circ S$ 的定义，必定存在中间元素（取 $y = x$）使得 $(x, x) \in R \circ S$ 成立。因此 $R \circ S$ 也是自反的，选项 D 正确。（注：传递性、反对称性、对称性在复合运算下一般不封闭，可轻易举出反例）。\\
	\textbf{故选 D。}
	
	\vspace{1em}
	
	\item \textbf{解:}\\
	根据给定的有向图，忽略自反的自环，图中的有向边（即关系）包含：$$E = \{ (3,1), (1,2), (3,4), (4,2), (3,2), (1,4) \}$$求其哈斯图需要求出覆盖关系，即剔除所有可通过传递性推导出的边。检查发现：
	\begin{enumerate}
		\item 由 $(3,1)$ 和 $(1,4)$，可传递推导出 $(3,4)$，故图中存在的 $(3,4)$ 应剔除。\\
		\item 由 $(1,4)$ 和 $(4,2)$，可传递推导出 $(1,2)$，故图中存在的 $(1,2)$ 应剔除。\\
		\item 由 $(3,1)$、$(1,4)$ 和 $(4,2)$（即 $3 \prec 1 \prec 4 \prec 2$），可推导出 $(3,2)$，故图中存在的 $(3,2)$ 应剔除。\\
	\end{enumerate}
	剔除多余的传递边后，剩下的直接覆盖关系为：$(3,1)$、$(1,4)$、$(4,2)$。这构成了一个全序关系（链）：$3 \prec 1 \prec 4 \prec 2$。反映在哈斯图上，应当是一条从下到上的单一直线，最底层是 3，往上依次是 1、4，最高层是 2。观察选项，只有 D 选项符合此拓扑结构。\\
	\textbf{故选 D。}
	
	\item \textbf{解:}\\
	已知哈斯图对应的全集 $A = \{1, 2, 3, 4, 5\}$，考查其子集 $B = \{2, 3, 4\}$。
	\begin{itemize}
		\item \textbf{最大元}：$B$ 中不存在大于等于 $B$ 中所有其他元素的元素（2 和 3 不可比），故\textbf{无}最大元。
		\item \textbf{最小元}：$B$ 中元素 4 满足 $4 \le 2$ 且 $4 \le 3$，故最小元为 \textbf{4}。
		\item \textbf{极大元}：$B$ 中没有严格大于 2 或 3 的元素，故极大元为 \textbf{2, 3}。
		\item \textbf{极小元}：$B$ 中没有严格小于 4 的元素，故极小元为 \textbf{4}。
		\item \textbf{上界}：在全集 $A$ 中，大于等于 $B$ 中所有元素（2, 3, 4）的元素仅有 \textbf{1}。
		\item \textbf{上确界}：所有上界中的最小者，显然为 \textbf{1}。
		\item \textbf{下界}：在全集 $A$ 中，小于等于 $B$ 中所有元素（2, 3, 4）的元素仅有 \textbf{4}（节点 5 不满足小于等于 2 和 4）。
		\item \textbf{下确界}：所有下界中的最大者，显然为 \textbf{4}。
	\end{itemize}
	将上述结果依次排列为：无, 4, 2, 3, 4, 1, 1, 4, 4。对比各项，与 D 选项完全一致。\\
	\textbf{故选 D。}
	
\end{enumerate}

\subsubsection{填空题}

\begin{enumerate}
	\item \textbf{解：}\\
	已知集合 $A = \{1, 2, 3, 4, 5, 6\}$, $B = \{1, 2, 3\}$，关系 $R = \{\langle x, y \rangle \mid x = y^2\}$ 且 $x \in A, y \in B$。\\
	当 $y = 1$ 时，$x = 1^2 = 1 \in A$，故 $\langle 1, 1 \rangle \in R$；\\
	当 $y = 2$ 时，$x = 2^2 = 4 \in A$，故 $\langle 4, 2 \rangle \in R$；\\
	当 $y = 3$ 时，$x = 3^2 = 9 \notin A$，不属于集合 A。\\
	因此，$R = \{\langle 1, 1 \rangle, \langle 4, 2 \rangle\}$。\\
	其逆关系 $R^{-1} = \{\langle y, x \rangle \mid \langle x, y \rangle \in R\} = \{\langle 1, 1 \rangle, \langle 2, 4 \rangle\}$。
	
	\vspace{1em}
	
	\item \textbf{解:}\\
	已知集合 $A = \{1, 2, 3, 4, 5, 6\}$，$R$ 是 $A$ 上的整除关系，即 $R = \{\langle x, y \rangle \mid x \text{ 整除 } y \}$。\\
	穷举符合条件的序偶：\\
	$1$ 可以整除 $1, 2, 3, 4, 5, 6$；\\
	$2$ 可以整除 $2, 4, 6$；\\
	$3$ 可以整除 $3, 6$；\\
	$4, 5, 6$ 分别只能整除自身。\\
	故 $R = \{\langle 1, 1 \rangle, \langle 1, 2 \rangle, \langle 1, 3 \rangle, \langle 1, 4 \rangle, \langle 1, 5 \rangle, \langle 1, 6 \rangle, \langle 2, 2 \rangle, \langle 2, 4 \rangle, \langle 2, 6 \rangle, \langle 3, 3 \rangle, \langle 3, 6 \rangle, \langle 4, 4 \rangle, \langle 5, 5 \rangle, \langle 6, 6 \rangle\}$。
	
	\vspace{1em}
	
	\item \textbf{解；}\\
	已知集合 $A = \{1, 2, 3, 4, 5, 6\}$, $B = \{1, 2, 3\}$，关系 $R = \{\langle x, y \rangle \mid x = 2y\}$ 且 $x \in A, y \in B$。\\
	当 $y = 1$ 时，$x = 2 \in A$，故 $\langle 2, 1 \rangle \in R$；\\
	当 $y = 2$ 时，$x = 4 \in A$，故 $\langle 4, 2 \rangle \in R$；\\
	当 $y = 3$ 时，$x = 6 \in A$，故 $\langle 6, 3 \rangle \in R$。\\
	因此，$R = \{\langle 2, 1 \rangle, \langle 4, 2 \rangle, \langle 6, 3 \rangle\}$。\\
	其逆关系 $R^{-1} = \{\langle 1, 2 \rangle, \langle 2, 4 \rangle, \langle 3, 6 \rangle\}$。
	
	\vspace{1em}
	
	\item \textbf{解:}\\
	已知集合 $A = \{a, b, c\}$ 上的二元关系 $R = \{\langle a, a \rangle, \langle a, b \rangle, \langle a, c \rangle, \langle c, c \rangle\}$。\\
	对称闭包 $s(R) = R \cup R^{-1}$。\\
	$R^{-1} = \{\langle a, a \rangle, \langle b, a \rangle, \langle c, a \rangle, \langle c, c \rangle\}$。\\
	故 $s(R) = \{\langle a, a \rangle, \langle a, b \rangle, \langle b, a \rangle, \langle a, c \rangle, \langle c, a \rangle, \langle c, c \rangle\}$。
	
	\vspace{1em}
	
	\item \textbf{解:}\\
	根据哈斯图，偏序关系 $R$ 需要满足自反性、反对称性和传递性。\\
	自反性：包含所有节点的自环，$\langle a, a \rangle, \langle b, b \rangle, \langle c, c \rangle, \langle d, d \rangle$。\\
	覆盖关系（图中直接相连的边，由下往上）：$\langle a, b \rangle, \langle a, c \rangle, \langle b, d \rangle, \langle c, d \rangle$。\\
	传递性：由 $a \prec b$ 且 $b \prec d$ 推导出 $\langle a, d \rangle$。\\
	综合可得：$R = \{\langle a, a \rangle, \langle b, b \rangle, \langle c, c \rangle, \langle d, d \rangle, \langle a, b \rangle, \langle a, c \rangle, \langle b, d \rangle, \langle c, d \rangle, \langle a, d \rangle\}$。
	
	\vspace{1em}
	
	\item \textbf{解:}\\
	集合 A上的偏序关系必须满足的三个性质是：\textbf{自反性}、\textbf{反对称性}、\textbf{传递性}。
	
	\vspace{1em}
	
	\item \textbf{解：}\\
	根据给定的关系图，写出关系 $R$ 的集合表示：\\
	$R = \{\langle 1, 2 \rangle, \langle 2, 1 \rangle, \langle 2, 3 \rangle, \langle 3, 4 \rangle\}$。\\
	根据复合关系的定义求 $R^2 = R \circ R$：\\
	因为 $\langle 1, 2 \rangle \in R, \langle 2, 1 \rangle \in R$，故 $\langle 1, 1 \rangle \in R^2$；\\
	因为 $\langle 1, 2 \rangle \in R, \langle 2, 3 \rangle \in R$，故 $\langle 1, 3 \rangle \in R^2$；\\
	因为 $\langle 2, 1 \rangle \in R, \langle 1, 2 \rangle \in R$，故 $\langle 2, 2 \rangle \in R^2$；\\
	因为 $\langle 2, 3 \rangle \in R, \langle 3, 4 \rangle \in R$，故 $\langle 2, 4 \rangle \in R^2$。\\
	所以 $R^2 = \{\langle 1, 1 \rangle, \langle 1, 3 \rangle, \langle 2, 2 \rangle, \langle 2, 4 \rangle\}$。
	
	\vspace{1em}
	
	\item \textbf{解：}\\
	既是等价关系又是偏序关系的例子是：\textbf{集合 A 上的恒等关系}（或写为 $I_A = \{\langle x, x \rangle \mid x \in A\}$）。\\
	（解释：等价关系要求对称性，偏序关系要求反对称性。若同时满足，则对于任意 $\langle x, y \rangle \in R$，必有 $\langle y, x \rangle \in R$，进而推出 $x = y$。再加上等价和偏序都要求的自反性，必然只有恒等关系满足条件。）
	
	\vspace{1em}
	
	\item \textbf{解:}\\
	已知集合 $A = \{1, 2, 3, 4, 5, 6\}$, $B = \{1, 2, 3\}$。\\
	关系 $R = \{\langle x, y \rangle \mid x = y^2, x \in A, y \in B\} = \{\langle 1, 1 \rangle, \langle 4, 2 \rangle\}$。\\
	$R$ 的关系矩阵 $M_R$（行对应 $A$ 中元素 $1\sim 6$，列对应 $B$ 中元素 $1\sim 3$）为：
	$$
	M_R = \begin{bmatrix}
		1 & 0 & 0 \\
		0 & 0 & 0 \\
		0 & 0 & 0 \\
		0 & 1 & 0 \\
		0 & 0 & 0 \\
		0 & 0 & 0
	\end{bmatrix}
	$$
	逆关系 $R^{-1} = \{\langle 1, 1 \rangle, \langle 2, 4 \rangle\}$。\\
	$R^{-1}$ 的关系矩阵 $M_{R^{-1}}$ 也就是 $M_R$ 的转置矩阵 $M_R^T$：
	$$
	M_{R^{-1}} = \begin{bmatrix}
		1 & 0 & 0 & 0 & 0 & 0 \\
		0 & 0 & 0 & 1 & 0 & 0 \\
		0 & 0 & 0 & 0 & 0 & 0
	\end{bmatrix}
	$$
	
	\vspace{1em}
	
	\item \textbf{解:}\\
	集合 A 上的等价关系的三个性质是：\textbf{自反性}、\textbf{对称性}、\textbf{传递性}。
	
	\vspace{1em}
	
	\item \textbf{解:}\\
	已知集合 $S = \{1, 2, 3, 4\}$ 上的关系 $R = \{\langle 1, 2 \rangle, \langle 2, 1 \rangle, \langle 2, 3 \rangle, \langle 3, 4 \rangle\}$。\\
	（1）求 $R \circ R$（也就是 $R^2$）：\\
	这与第 7 题的集合对应关系完全一致，由复合关系定义可得：\\
	$R \circ R = \{\langle 1, 1 \rangle, \langle 1, 3 \rangle, \langle 2, 2 \rangle, \langle 2, 4 \rangle\}$。\\
	（2）求 $R^{-1}$：\\
	将 $R$ 中所有序偶的元素位置互换即可得到逆关系：\\
	$R^{-1} = \{\langle 2, 1 \rangle, \langle 1, 2 \rangle, \langle 3, 2 \rangle, \langle 4, 3 \rangle\}$。
	
\end{enumerate}















